% concatenated from revised.tex


\documentclass[letterpaper, 10 pt, conference]{ieeeconf}
\IEEEoverridecommandlockouts
% The preceding line is only needed to identify funding in the first footnote. If that is unneeded, please comment it out.
\let\proof\relax
\let\endproof\relax
\let\labelindent\relax
\overrideIEEEmargins  
\setlength{\topmargin}{-10pt}
\usepackage{comment}
\usepackage{booktabs}
\usepackage{cite}
\usepackage{mathtools}
\usepackage{amsmath,amssymb,amsfonts}
\usepackage{algorithmic}
\usepackage{graphicx}
\usepackage{textcomp}
\usepackage{makecell}
\usepackage{xcolor}
\usepackage{amsthm}
\usepackage{microtype}
\usepackage{enumitem}
\newtheorem{definition}{Definition}
\newcommand{\N}{\mathbb{N}}
\newcommand{\Z}{\mathbb{Z}}
\newcommand{\Q}{\mathbb{Q}}
\newcommand{\R}{\mathbb{R}}
\newcommand{\C}{\mathbb{C}}
\newcommand{\E}{\mathbb{E}}
\newcommand{\Var}{\mathrm{Var}}
\newcommand{\Cov}{\mathrm{Cov}}
\newcommand{\Prob}{\mathbb{P}}
% \newcommand{\abs}[1]{\lvert#1\rvert}
% \newcommand{\norm}[1]{\lVert#1\rVert}
\newcommand{\set}[1]{\{#1\}}
\newcommand{\ceil}[1]{\lceil#1\rceil}
\newcommand{\floor}[1]{\lfloor#1\rfloor}
\newcommand{\argmax}{\operatorname*{argmax}}
\newcommand{\argmin}{\operatorname*{argmin}}
\newcommand{\ip}[2]{\left\langle #1,#2\right\rangle}
\usepackage{amsmath}
\DeclareMathOperator*{\esssup}{ess\,sup}

\newcommand{\mcl}[1]{\mathcal{#1}}
\newcommand{\mbb}[1]{\mathbb{#1}}
\newcommand{\mbf}[1]{\mathbf{#1}}
\newcommand{\mrm}[1]{\mathrm{#1}}
\newcommand{\norm}[1]{\left\lVert #1 \right\rVert}

\newcommand{\mat}[1]{\begin{matrix} #1 \end{matrix}}
\newcommand{\pmat}[1]{\begin{pmatrix} #1 \end{pmatrix}}
\newcommand{\vmat}[1]{\begin{vmatrix} #1 \end{vmatrix}}
\newcommand{\smat}[1]{\begin{smatrix} #1 \end{smatrix}}
\newcommand{\bmat}[1]{\begin{bmatrix} #1 \end{bmatrix}}
\newtheorem{theorem}{Theorem} % Can define other types too
\newtheorem{corollary}{Corollary}
\theoremstyle{definition}
\newtheorem{example}{Example}
\def\BibTeX{{\rm B\kern-.05em{\sc i\kern-.025em b}\kern-.08em
    T\kern-.1667em\lower.7ex\hbox{E}\kern-.125emX}}
\allowdisplaybreaks
\setlength{\abovedisplayskip}{1mm}   % equation Space above
\setlength{\belowdisplayskip}{1mm}   % equation Space below
\setlength{\abovecaptionskip}{0mm} % figure space
\setlength{\belowcaptionskip}{0mm} % figure space below
\raggedbottom

\begin{document}

\title{\LARGE \bf Impulse-to-Peak-Output Norm Optimal State-Feedback Control of Linear PDEs}%
\author{Tristan Thomas$^{1}$, Sachin Shivakumar$^{2}$, and Javad Mohammadpour Velni$^{1}$
\thanks{$^{1}$Tristan Thomas \{{\tt\small tthoma7@clemson.edu}\} and Javad Mohammadpour Velni \{{\tt\small javadm@clemson.edu}\} are with the Department of Mechanical Engineering at Clemson University.}%
\thanks{$^{2}$Sachin Shivakumar \{{\tt\small sshivakumar@lanl.gov}\} is with the Center for Nonlinear Studies, Los Alamos National Laboratory. }%
\thanks{This work is partly supported by the LDRD program (20250614CR-NLS) of Los Alamos National Laboratory. LA-UR-26-22329. Tristan Thomas is supported by the U.S. Department of Education through GAANN program P200A220075.}
}


%\author{\authorblockN{Tristan Thomas}
%\authorblockA{\textit{Department of Mechanical Engineering} \\
%\textit{Clemson University}\\
%Clemson, SC \\
%tthoma7@clemson.edu}
%\and
%\authorblockN{Sachin Shivakumar}
%\authorblockA{\textit{Center for Nonlinear Studies} \\
%\textit{Los Alamos National Laboratory}\\
%Los Alamos, NM \\
%sshivakumar@lanl.gov}
%\and
%\authorblockN{Javad Mohammadpour Velni}
%\authorblockA{\textit{Department of Mechanical Engineering} \\
%\textit{Clemson University}\\
%Clemson, SC \\
%javadm@clemson.edu}
%}

\maketitle

\begin{abstract}
Impulse-to-peak response (I2P) analysis for state-space ordinary differential equation (ODE) systems is a well-studied classical problem. However, the computational techniques employed for I2P optimal control of ODEs have not been extended to partial differential equation (PDE) systems due to the lack of a computation-friendly universal transfer function and state-space representation. Recently, however, partial integral equation (PIE) representation was proposed as the desired state-space representation of a PDE, and Lyapunov stability theory was used to solve various control problems, such as stability and $H_\infty$-optimal control. In this work, we utilize this PIE framework, and associated Lyapunov techniques, to formulate the I2P response analysis problem as a solvable convex optimization and obtain provable bounds for the I2P-norm of linear PDEs. Moreover, by establishing strong duality between primal and dual formulations of the optimization problem, we develop a constructive method for I2P optimal state-feedback control of PDEs and demonstrate the effectiveness of the method on various examples.



% Analysis of peak output response to impulse input is crucial in dynamical systems, especially in physical systems where a momentary deviation from prescribed limits can compromise structural integrity and cause catastrophic failures. Despite the importance, analyzing impulse to peak output response for infinite dimensional systems, such as partial differential equations (PDE), can be difficult because impulse inputs and peak outputs are typically parametrized by operators that are unbounded in time. We circumvent this difficulty by using a Lyapunov function based characterization of the impulse-to-peak norm and by using the partial integral equation (PIE) representation of PDEs to derive a convex optimization approach to establish provable bounds on impulse-to-peak-output norm. 

% We show that the resulting optimization problem has Partial Integral (PI) operators as decision variables and involves conic constraints on PI operators---i.e., a linear PI inequality (LPI) problem. Moreover, we show that the PIE representation itself admits a dual PIE representation with same impulse-to-peak-output norm, resulting in a complementary dual LPI formulation for the norm bounding problem. Using this dual LPI, we derive a convex synthesis conditions to design impulse-to-peak-output optimal state-feedback controller for linear PDE systems and demonstrate the application on various numerical examples.

%This work presents results on computation of the optimal impulse-to-peak system norm for partial differential equation (PDE) systems. This method exploits the partial integral equation (PIE) framework to represent a given PDE as a PIE, and formulate the computation of the defined system norm as a convex optimization problem using that PIE representation. We present results that show the development of the linear PI inequality (LPI) optimization problem we solve for obtaining the norm, and also prove that the norm of the primal and dual norms are equivalent both in theory and in simulation.
\end{abstract}

%\begin{IEEEkeywords}
%Convex Optimization, Partial Differential Equation, Optimal Control
%\end{IEEEkeywords}

\section{Introduction}
%\textit{PDE models are widely used $\to$ thus, optimal control is well-studied $\to$ optimal control does not focus on important transient behavior $\to$ example: battery $\to$ instead of optimal control, adaptive methods are used which are computationally expensive $\to$ We need offline computation like in ODEs $\to$ there are no such methods for PDEs because {next paragraph}.} 

Optimal control of partial differential equations (PDEs) has been a well-studied topic \cite{Moura,vanKeulen,troltzsch2010optimal} owing to their ubiquitous use in a range of practical applications. However, metrics used therein cannot certify safety-critical transient behaviors. For instance, instantaneous spikes in current or temperatures that often lead to failure/degradation in Lithium-ion batteries are not captured by $H_\infty$/LQR metrics~\cite{mouraCST}. Hence, one often uses optimal control approaches such as MPC \cite{Guay} to enforce such transient specifications. Although widely used in ordinary differential equation (ODE) systems, such predictive methods require forward simulation of the system dynamics -- a computationally expensive step for PDEs. Since the goal is to design a computationally inexpensive control that can certify transient behavior, unlike $H_\infty$/LQR-optimal control, impulse-to-peak (I2P) norm optimal control is a natural candidate. Hence, we develop a computational method to design I2P optimal state-feedback controllers for a large class of PDEs.  

%Many engineered, natural, and social systems  are modeled using partial differential equations (PDEs)\cite{icdse25}. For controlling PDE systems, both non-optimal\cite{VAZQUEZ2026112572} and optimal (LQR/LQG \cite{Moura}, $H_2/H_\infty$ \cite{vanKeulen}) control methods have been developed. However, commonly used methods for optimal control of PDEs are ill-suited for certain safety-critical applications, where certified transient performance of metrics like rise time or peak response is crucial. For example, in lithium-ion batteries, short duration high amplitude spikes in current or ambient temperature can push internal cell temperatures above critical thresholds---a common cause of failures and battery degradation \cite{mouraCST}. This motivates adaptive methods, such as MPC \cite{Guay} and control barrier function (CBF) based methods \cite{wang2023safe}, which may provide transient performance or safety guarantees. These methods, however, may lack general applicability to different PDEs and can be computationally expensive. Thus, there is need to consider offline optimal control methods that are better suited for transient performance guarantees.

%In this work, we will develop a computational method to design provably optimal controllers for a large class of linear partial differential equation (PDE) systems that minimize impulse-to-peak norm. Unlike for ordinary differential equation (ODE) systems, where this problem is well-studied (e.g., \cite{TOKUNAGA19981473}), there is a lack of computational methods to obtain such controllers for infinite dimensional PDE systems. %, where computationally efficient techniques for provable bounds and controller synthesis are lacking. 
%Commonly used optimal control methods for PDEs, such as LQR/LQG \cite{Moura} or $H_2/H_\infty$-optimal control \cite{vanKeulen}, are ill-suited for certain safety-critical applications, where certified transient performance such as rise time or peak response is required. In systems like lithium-ion batteries, short duration high amplitude spikes in current or ambient temperature can push internal cell temperatures above critical thresholds---a common cause of failures and battery degradation \cite{mouraCST}. %discusses the severity of thermal faults in Lithium-ion batteries, which are responsible for a significant portion of failures and degradation. Short, high-amplitude disturbances like current surge or ambient temperature spikes can push internal cell temperatures above critical thresholds.
%Due to the inability of classical optimal control methods to address transient behavior problems, instead, one must rely on adaptive controls, such as MPC \cite{Guay} and CBF \cite{wang2023safe}. However, these methods often require demanding online computation or lack generality/optimality.

\begin{comment}
The importance of general peak norm bounding, including impulse-to-peak norm, for dynamical systems lies in ensuring safety and understanding worst case response of these systems. As a result, general methods of peak bounding for various system classes and various applications have been developed. In \cite{MILLER2026112641}, the authors formulate an infinite-dimensional measure program along with its semidefinite program (SDP) relaxation for estimating the peak response of hybrid systems. Impulse-to-peak optimization has been explored using dc programming in \cite{ZHAO20235059} for positive linear systems. Impulse-to-peak synthesis methods have found application in the control of an aerospace launcher discussed in \cite{1393878}. While there appears to be no work on applying impulse-to-peak synthesis methods for practical applications for PDE systems, it is the belief of the authors that this is a direct artifact of the lack of desirable methods for such a task.
\end{comment}
\enlargethispage{\baselineskip}
Due to simplicity and maturity of tools available for ODE optimal control, the standard approach for I2P optimal control of PDEs often involves approximating a PDE by an ODE (\emph{early-lumping}) and then using ODE methods, such as \cite{TOKUNAGA19981473}, to design an I2P-optimal control. However, approximation methods often introduce truncation errors---resulting in spillover effects and a lack of closed-loop stability or performance bounds \cite{balas}. \emph{Late-lumping} approaches, on the other hand, use variational calculus, Hamilton-Jacobi-Bellman (HJB) equation, or a Riccati PDE and can provide provable performance \cite{Lasiecka_Triggiani_2000}, but the controllers must be obtained by solving nonlinear PDEs with unbounded input operators, which cannot be solved exactly in all but the simplest systems. While \cite{mironchenko2020input} developed Lyapunov characterization of this problem, and sidesteps the above limitations, it does not provide a constructive method to find controllers. Thus, existing approaches typically lack \textit{provable} performance, \textit{generality}, or both.  

%While traditional offline optimal control methods for PDEs avoid computational burden, they too often lack generality and do not possess provable optimality. These approaches can largely be categorized as either ``early-lumping'' or ``late-lumping'' approaches. Early lumping approaches typically approximate PDEs by ODEs either by discretization or projection. The choice of discretization schemes and projection bases often depend on the specific PDE to ensure numerical stability and capture important physics.
%The discretization schemes differ for different PDEs, e.g., finite-difference, finite-volume, and other methods. Projection based methods for early-lumping require a choice of basis. Due to the difference in spectral properties of the spatial operator of a PDE as well as the regularity of the solutions, different choices and sizes of bases are required. 
%Thus, a single choice of early-lumping approach does not generalize for all PDEs. Moreover, due to spillover effects \cite{balas}, there is no guarantee of closed-loop stability or performance bounds. Late-lumping approaches typically use variational calculus, the Hamilton-Jacobi-Bellman (HJB) equation, or  a Riccati PDE based solution to find the optimal control. However, the resulting optimal controller often does not have a closed-form solution (or may not even be constructive) and require numerical approximations, which are difficult to solve, particularly, when input operators are unbounded. In general, the late-lumping methods guarantee stability and performance only as a convergence limit.

%In this work, we circumvent issues with adaptive and traditional offline optimal control methods to \textit{synthesize controllers offline} that can achieve \textit{provable transient performance}. In this case, the transient metric of interest is the \textit{impulse-to-peak norm}. By minimizing this norm, we can reduce a system's peak response to pulse disturbances, enhancing safety and reliability in systems such as the aforementioned lithium-ion batteries.

Naturally, to develop a general algorithm that computes controllers with provable performance, one must first separate the method from the model. A traditional example of this is linear matrix inequality (LMI)-based methods for ODE analysis/control, where a single synthesis method works for all ODE systems of a particular class. Unfortunately, a similar method is not available for PDEs due to lack of a general representation that admits such finite-time algorithms. While the semigroup approach \cite{curtain2012introduction} does provide a general state-space representation for PDEs, it is not convenient for LMI-style computational methods. To overcome this, we use the partial integral equation (PIE) representation \cite{shivakumar_extension_2024}---an \textit{equivalent} state-space-like representation for linear PDEs parameterized by bounded integral operators called partial integral (PI) operators. Moreover, PI operators are closed under composition, addition, and adjoint. Thus, unlike with PDEs, various LMI-style conditions with PI operator variables can analogously be defined and \emph{solved} for PIEs. %means inequalities involving these operations are well defined fully in the space of PI operators. The algebraic closure, combined with the fact that PI operators can be parameterized as finite-dimensional matrices, allows us to represent PI operator constraints in tractable semidefinite programs (SDPs). 


\begin{comment}
Using the partial integral equation (PIE) framework for representing PDEs, general methods for other norm bounding problems for PDEs have been developed \cite{shivakumar_dual_2024}. This work takes a similar approach exploiting the PIE representation framework to develop a convex optimization-based approach for the impulse-to-peak problem. The major benefit in doing so is that the method is PDE agnostic. It is designed for any linear PDE which admits a PIE representation. Additionally, PIEs are parameterized by partial integral (PI) operators, which form a closed algebra. This algebraic structure allows us to compute numerically the composition and addition of PI operators, developing synthesis methods entirely in the operator space. Finally, PI operators are bounded integral operators, which can be parameterized by matrices. This allows us to optimize over the space of positive integral operators by optimizing over the space of positive definite matrices. This detail is discussed further in section II.C.
\end{comment}
This framework of using PIEs to represent PDEs, posing PDE optimal control problem as an equivalent PIE optimal control problem, and using associated computational machinery to solve the PIE control problems has been successfully demonstrated for $H_\infty$-optimal control problems of PDEs \cite{shivakumar_dual_2024}. In this work, we will leverage the universality and computational benefits afforded by the PIE framework to solve I2P-optimal control problem. In that regard, the \textit{key contributions} of this work are: 1) characterizing the I2P-norm of a PDE in terms of its PIE representation, 2) proving equivalence of the I2P-norm for a PIE and its dual, 3) posing the open-loop I2P-norm upper bounding problem and the closed-loop I2P-norm minimization problem as \emph{solvable} convex optimization problems. The resulting analysis and synthesis conditions are computationally implemented using PIETOOLS \cite{shivakumar_pietools_2025} and have been demonstrated on various numerical examples.

%%%%%%%%%%%%%%%%%%%%%%%%%%%%%%%%%%%%%%%%%%%%%%%%%%%%%%%%%%%%%%%%%%%%%%%

\section{Notation and Preliminaries}
We will use lowercase font for functions of time and space, e.g., $x(t)$ and $\mbf x(t,s)$, and uppercase calligraphic font $\mcl A$ for operators on functions. Given $\mcl A$, an operator on a Hilbert space, we denote the adjoint operator w.r.t. canonical inner-product by $\mcl A^*$.
%, i.e., $\ip{x}{\mcl A y} = \ip{\mcl A^*x}{y}$ for all $x,y$. 
Partial derivatives w.r.t. time and space are denoted by $\partial_t$ and $\partial_s$.
$L_2^n[X]$ denotes the space of $\R^n$-valued Lebesgue square-integrable functions on $X$ and is equipped with the canonical inner-product $\ip{\cdot}{\cdot}_{L_2^n[X]}$. For functions on $\R_+$, we define the supremum (peak) norm as $\norm{x}_\infty = \esssup_{t\ge 0} \norm{x(t)}$ and associated Banach space as $L_\infty[\R_+]$. %The set of bounded linear operators mapping between Hilbert spaces $X$ and $Y$ is denoted by $\mbf L(X,Y)$. For any $\mcl P\in \mbf L(X,Y)$, we denote the adjoint by $\mcl P^*\in \mbf L(Y,X)$. 
Lastly, we abbreviate $\R^m\times L_2^n[X]$ as $\R L_2^{m,n}[X]$, e.g., if $\begin{bmatrix}
    x \\ \mathbf{x}
\end{bmatrix} \in  \mathbb{R}L_2^{m,n}$ then $x \in \mathbb{R}^m$ and $\mathbf{x}\in L_2^n$.

Next, similar to \cite{shivakumar_dual_2024}, we briefly introduce the Partial Integral (PI) operators, Partial Integral Equations (PIEs), and Linear PI Inequalities (LPIs) optimization problems.

\subsection{The Class of Partial Integral Operators}

4-PI operators, a subclass of general PI operators, are used for the parameterization of PIEs just as matrices are used for the parameterization of linear ODEs. 

\begin{definition}
Given $x \in \mathbb{R}^{m_1}$ and $\mathbf{x} \in L_2^{n_1}[a,b]$, we say $\mathcal{P}:\R L_2^{m_1,n_1}\to \R L_2^{m_2,n_2}$ is a \textbf{4-PI operator} if it can be parameterized as \begin{align}
\left(\mcl P \begin{bmatrix}
    x\\ \mathbf{x}
\end{bmatrix}\right)(s) \coloneqq \begin{bmatrix}
    Px + \int_a^b{Q_1(\theta)x(\theta) d \theta }\\
    Q_2(s)x + \mathcal{R}\mathbf {x}(s)\nonumber
\end{bmatrix},\end{align}
where
\begin{align}
    (\mathcal{R}\mathbf{x})(s)\!\!&=\!\!R_0(s)\mathbf{x}(s) \!+\!\! \int_a^s\!\!{R_1(s,\theta)\mbf{x}(\theta)d\theta}\!+\!\!\int_s^b\!\!{R_2(s,\theta)\mathbf{x}(\theta)d\theta}\nonumber
\end{align}
for some matrix $P$ and matrix-valued polynomials $Q_1, Q_2, R_0, R_1,$ and $R_2$. Given $P,Q_i,R_i$, we denote 4-PI operators in the compact form
$\mcl P= \left[ \begin{array}{c|c}
P & Q_1 \\
\hline
Q_2 & \{R_i\}
\end{array} \right]$.
We denote the class of 4-PI operators as $\Pi_4$. %\subset \mathcal{L}(\mathbb{R}L_2^{m_1,n_1},\mathbb{R}L_2^{m_2,n_2})$.
\end{definition}
% In the above definition, $\mathcal R$ itself is called a 3-PI operator and a set of such operators is denoted by $\Pi_3$.
%These operators (both 3-PI and 4-PI) are defined with $n_1, n_2$ not necessarily equal to 1, however, in control applications, typically $n_1 = n_2 = 1$. Additionally, 
The parameter dimensions are inherited from the dimensions of the map, i.e., $P \in \mathbb{R}^{m_2 \times m_1}$, $Q_1(s) \in \mathbb{R}^{m_2 \times n_1}$, $Q_2(s)\in \mathbb{R}^{n_2\times m_1}$, and $R_0(s), R_1(s,\theta), R_2(s,\theta) \in \mathbb{R}^{n_2 \times n_1}$. In the special case $(m_1,n_1) = (m_2,n_2)$, $\Pi_4$ forms a $*$-subalgebra \cite{shivakumar_extension_2024}. Notably, 4-PI operators are always bounded since their kernels are polynomials.  

\subsection{The Class of Partial Integral Equations}
Next, we examine the class of dynamical systems that are parametrized by these 4-PI operators.
\begin{definition}
    Given 4-PI operators $\{\mcl T, \mcl A, \mcl B, \mcl C, \mcl D\}\subset \Pi_4$, we say $\mbf x(t,\cdot)\in \R L_2^{m,n}, z:\R_+ \to \R^{n_z}$ are governed by a {\bf Partial Integral Equation (PIE)} if $\{\mbf x,z\}$ satisfies
    \begin{align}
        \label{eq:PIE}
        \partial_t\left(\mcl T\mbf x\right)(t) &= \mcl A \mbf x(t) \!+\! \mcl B w(t),\quad z(t) = \mcl C \mbf x(t)\!+\!\mcl D w(t),
    \end{align}
    for some initial conditions $\mcl T\mbf x(0)$ and $w(t) \in \mbb{R}^{n_w}$.
    %A dynamical system of the form Eq. \eqref{eq:PIE}.% is denoted by $G(\mcl T,\mcl A, \mcl B, \mcl C, \mcl D)$.
\end{definition}
The initial conditions for PIEs are often defined on $\mcl T\mbf x$ instead of $\mbf x$ because, by the nature of PIE representation of PDEs, initial conditions on $\mcl T\mbf x$ are wellposed but not necessarily on $\mbf x$.

Since the 4-PI operators are linear, bounded, and form a *-subalgebra, they are a natural extension of state-space ODE representation to infinite-dimensional spaces. In fact, various linear PDE and time-delay systems \cite{shivakumar_extension_2024} can be described using the form Eq. \eqref{eq:PIE}. For example, consider a transport equation, $\partial_t v(t,s) = \partial_s v(t,s)$ with boundary condition $v(t,0) = 0$. By defining $\mbf{x}(t,s) = \partial_s v(t,s)$ and using the Fundamental Theorem of Calculus, $v(t,s) = \int_0^s \mbf{x}(t,\zeta) \ d\zeta$. This is of the form $\partial_t(\mcl T \mbf{x})(t) = \mcl A\mbf{x}(t)$ where $\mcl T\mbf{x} = \int_0^s \mbf{x}(t,\zeta) \ d\zeta $, and $\mcl A = 1$. We illustrate the class of PDEs considered in this work in the following subsection.


%we get $\int_0^s \mbf x(t, \zeta)\ d\zeta = v(t,s)-v(t,0)$, which reduces to

\subsection{The Class of Partial Differential Equations}
Having presented a general class of PIEs, we next present the class of PDEs that admit such a representation. %While the full class of PDEs that admit a PIE representation is beyond the scope of this discussion (see \cite{shivakumar_extension_2024} for details),
For illustration, we consider a class of $2^{nd}$-order PDEs in a single spatial variable on a compact domain, $s \in [a, b]$. 
Given a PDE state $\mbf x = [\mbf x_1, \mbf x_2,\mbf x_3]^\top$, where spatial and time dependence are sometimes omitted for brevity, we can represent such PDEs as 
\begin{align}\label{eq:PDE}
    \partial_t \mbf{x}(s,t) &= \mcl A_p \mbf x(s,t) + B_{21}(s)w(t) + B_{22}(s)u(t),\\
    z(t) &= \mcl C \mbf x(t) + \mcl D_{11} w(t) + \mcl D_{12} u(t),\nonumber
    \end{align}
    $\mbf x_c(s,t) = 
        [\mbf x_2, \mbf x_3, \partial_s\mbf x_{3}]^\top$,
where the operators $\mcl A_p$ and $\mcl C$ are defined as
\begin{align}
    (\mcl A_p \mbf x)(s,t)
    &\coloneqq A_0(s) \mbf x+ A_1(s)[\partial_s\mbf x_{2},
        \partial_s\mbf x_{3}]^\top+A_2(s) \partial_s^2\mbf x_{3} \nonumber,\\
    (\mcl C \mbf x)(t) &\coloneqq \begin{multlined}[t]
         C_{10} [
        \mbf x_c(a,t),\mbf x_c(b,t)]^\top + \int_a^b \Bigl(C_a(s)\mbf{x}\\+C_b(s) 
        [\partial_s\mbf x_{2},\partial_s\mbf x_{3}]^\top\Bigr) ds.  \end{multlined}\nonumber
\end{align}
For all $t\ge 0$, $\mbf x(\cdot, t)$ lies in the domain
   \[D(\mcl A_p) \!\coloneqq\! \left\lbrace \mat{\mbf{x}(s,t) \in L_2^{n_1}[a,b] \times W_{2,1}^{n_2} [a,b] \times W_{2,2}^{n_3}[a,b] \colon \nonumber \\B
        [\mathbf{x}_c(a), \mathbf{x}_c(b)]^\top = 0} \nonumber\right\rbrace,\]
where $B$ defines boundary conditions on the $\mbf x$, and $W_{2,p}^n$ is the canonical Sobolev space, i.e., $\mbf x\in W_{2,p}^n$ if $\partial_s^k\mbf x\in L_2^n$ for all $k\le p$.
If the conditions of \textit{PIE-compatibility} (see \cite[Sec. IV-A]{shivakumar_extension_2024}) are satisfied, any PDE parameterized in the above form can be equivalently represented as a PIE \eqref{eq:PIE}.



\subsection{Solving LPIs and 4-PI Operator Inversion}
Any \textit{PIE-compatible} PDE of the form \eqref{eq:PDE} can be represented as PIEs using explicit formulae in \cite{shivakumar_extension_2024} (or by using the open-source PIETOOLS toolbox \cite{shivakumar_pietools_2025}). Thus, we implicitly assume that the PDE is PIE compatible and the PIE parameters, ($\mcl{T}, \mcl{A}, \mcl{B}, \mcl{C}, \mcl{D}$), are obtained using available tools. In addition, the results herein are presented as linear PI inequality (LPI) optimization problems, which have the form
\begin{align}\label{eq:LPI}
\min_{\mcl P\in \Pi_4}&\quad   f(\mcl P),\quad 
s.t., \quad g_i(\mcl P) = 0,\; h_j(\mcl P)\le 0,
\end{align}
where $f$, $g_i$, and $h_j$ are linear functionals involving PI decision variables $\mcl P$. For example, given a PIE of the form $\partial_t (\mcl T\mbf x)(t) = \mcl A \mbf x(t)$, Lyapunov stability test can be formulated as
\[
\mcl P\succ 0, \;\; \mcl T^*\mcl P\mcl A+\mcl A^*\mcl P\mcl T\preceq 0.
\]
Problems of the above form are solved computationally by parameterizing positive operators using positive matrices. Specifically, one can parameterize $\mcl P = \mcl Z^*P\mcl Z$ using matrix $P>0$ and constrain
\[
\mcl T^*\mcl Z^*P\mcl Z\mcl A+ \mcl A^*\mcl Z^*P\mcl Z\mcl T = -\mcl Z_d^* Q\mcl Z_d,
\]
where $\mcl Z$ and $\mcl Z_d$ are 4-PI operator bases and matrix $Q\ge 0$. Since, 4-PI operator kernels are polynomials, the above constraint is merely a polynomial equality constraint, and thus, is equivalent to constraints on coefficients of polynomials stored in $P$ and $Q$. Thus, infinite-dimensional LPI problem is transformed to finite-dimensional LMI problem and can be solved using interior-point methods. Details of the implementation and procedure of solving LPI constrained problems can be found in \cite{shivakumar_pietools_2025} and \cite{PEET2021109473}.
%Note that the choice of a basis does not imply truncation or approximation of the PDE. It simply indicates a restriction for the search of a particular 4-PI operator, $\mcl P$, that makes the above operator equation true. The PIE in this case is parameterized only by $\mcl T, \mcl A$, both of which are not approximated using a basis.


\begin{comment}
To illustrate, we show how PI operators are parameterized by matrices with an illustrating example. For this, we restrict Definition 1 to the class of 3-PI operators. This is simply a 4-PI operator that does not act on an ODE state. Thus, the 3-PI operator, $\mcl M$ can be written as $\mcl M = \mcl R \mbf x(s) =  R_0(s) \mbf{x}(s) + \int_a^s R_1(s,\theta) \mathbf{x}(\theta) d\theta + \int_s^b R_2(s,\theta) \mbf{x}(\theta) d\theta $. We know that we can represent $\mcl M = Z^* P Z$ for $P>0$ as mentioned previously. If we consider a basis for the kernels as a monomial basis of degree 0 and a 1D PDE for simplicity we will obtain $\mcl Z = \begin{bmatrix}
    1 & \int_a^s (\cdot)d\theta & \int_s^b(\cdot)d\theta
\end{bmatrix}^T$ and $P = \begin{bmatrix}
    P_{11}&P_{12}&P_{13}\\ 
    P_{12}&P_{22}&P_{23}\\
    P_{13}&P_{23}&P_{33}
\end{bmatrix}$. The total product $\mcl Z^* P\mcl Z$ in this case would yield, $P_{11}$
\end{comment}

Lastly, controller synthesis problems often require inversion of 4-PI operators, which can be calculated numerically using the formulae in \cite[Sec. VII]{shivakumar_dual_2024}. The PDE to PIE conversion, LPI optimization, and controller gain reconstruction will be performed using PIETOOLS.

%A detailed discussion on numerical nuances, such as conversion of PDEs to PIEs, setting up PI operator-valued optimization problems, and reconstructing controller gains by computing PI operator inverses, is beyond the scope of this paper and will be omitted. However, we note that all these steps can be performed computationally and all the results presented herein are obtained by employing the functionality offered by PIETOOLS.

\section{Duality in PIEs}

Every PIE-compatible PDE has an \textit{equivalent} PIE representation, i.e., well-posedness, stability, and $L_2$-gain bound of the PDE implies the same for its PIE representation. However, the properties of the dual system are more critical to controller synthesis problems. Unlike the dual of a PDE, the dual of a PIE is readily available and is defined below.

\begin{definition}\label{def:dualPIE}
Given a PIE Eq. \eqref{eq:PIE} defined by 4-PI operators $\{\mcl T, \mcl A,\mcl B, \mcl C,\mcl D\}\subset \Pi_4$, we say $\bar{\mbf x}:\R_+\to \R L_2^{m,n}$, $\bar z:\R_+\to \R^{n_w}$ are governed by \textbf{the dual PIE}, if $\{\bar{\mbf x},\bar z\}$ satisfies 
\begin{align}
\partial_t(\mcl T^*  \bar{\mbf x})(t)\!=\! \mcl A^* \bar{\mbf x}(t)\!+\!\mcl C^* \bar{w}(t),\; \bar{z}(t)\!=\!\mcl B^*  \bar{\mbf x}(t)\!+\!\mcl D^*\bar w(t), \label{PIEd}
%\qquad y(t)=\mcl C_2  \mbf x(t) + \mcl D_{21} w(t),
\end{align}
for some initial conditions $\mcl T^*\bar{\mbf x}(0)$ and $\bar w\in L_2^{n_z}[\R_+]$.
\end{definition}
    Since the PI parameters of the primal system, Eq. \eqref{eq:PIE}, lie in a *-subalgebra, the dual system, Eq. \eqref{PIEd}, is also parameterized by PI parameters and admits the same form. Moreover, one can show that a PIE and its dual possess the same stability property and $L_2$-gain bound \cite{shivakumar_dual_2024}.%---an attribute of PIE systems referred to as \textcolor{blue}{\textbf{duality}}.

Both the structural and behavioral equivalence between the primal and dual play a crucial role in the controller synthesis problems. The former allows any computational method developed for the primal to be applicable for the dual. The latter allows one to indirectly establish properties of the primal by proving properties of its dual. For example, the state-feedback controller synthesis for primal is difficult since the Lyapunov stability conditions for the closed-loop system are bilinear (non-convex) in Lyapunov function parameter ($\mcl P$) and controller gains ($\mcl K$): 
\[
\mcl P\succ 0, \; (\mcl A+\mcl B\mcl K)^*\mcl P\mcl T+\mcl T^*\mcl P(\mcl A+\mcl B\mcl K) \preceq 0.
\]
However, since the dual is defined by the same class of parameters and inherits the same stability properties, one can instead find a stabilizing controller by examining Lyapunov conditions for the \textbf{dual} of the closed-loop system, which is a convex problem in decision variables ($\mcl P, \mcl Z=\mcl K\mcl P$):
\[
\mcl P\succ 0, \; (\mcl A\mcl P+\mcl B\mcl Z)\mcl T^*+\mcl T(\mcl P\mcl A^*+\mcl Z^*\mcl B^*) \preceq 0.
\]
By solving the dual constraints, one can find the controller gains $\mcl K = \mcl Z\mcl P^{-1}$ and the associated Lyapunov function parameter $\mcl P$ that proves the closed-loop stability of the primal system. Our objective in this paper is to extend this method for impulse-to-peak norm optimal controller synthesis, and this requires:
\begin{itemize}
    \item finding optimization-based characterization of the impulse-to-peak norm of a PIE Eq. \eqref{eq:PIE}.
    \item showing that impulse-to-peak norm of the primal Eq. \eqref{eq:PIE} and its dual Eq. \eqref{PIEd} are equal.
\end{itemize}

% The dual representation presented above is crucial because convex formulations of the controller synthesis problems often use dual systems, similar to dual state-space ODE representations used in LMI-based methods. 

% Utility of the dual PIE representation for control design will be evident in the later sections. However, first, we must define the optimal impulse-to-peak output control problem and establish that the duality in PIEs holds true for this problem.


\subsection{Impulse-to-Peak Norm Calculation of a PIE}
First, we will formally define the impulse-to-peak norm of a PIE and establish the equivalence of this system norm in primal and dual PIEs. 
\begin{definition}\label{def:ip-normsystem}
    Given $\{\mcl T, \mcl A, \mcl B, \mcl C\}$, we use $G(\mcl T, \mcl A,\mcl B,\mcl C)$ to denote a PIE with impulse inputs of the form
\begin{align}
\partial_t (\mcl T \mbf x)(t)&= \mcl A \mbf x(t)\!+\!\mcl B w(t),\; \mcl T\mbf x(0) =0,\; z(t)=\mcl C \mbf x(t), \label{PIE2}
%\qquad y(t)=\mcl C_2  \mbf x(t) + \mcl D_{21} w(t),
\end{align}
where $\mbf x(t) \in \R L_2^{m,n}[a,b]$ is the state, $w(t)$ is an \textbf{impulsive disturbance} of the form $w(t) = \delta(t)v$ with $v \in \mathbb{R}^{n_w}$, and $z(t) \in \R^{n_z}$ is the output.
\end{definition} Then, we can define the impulse-to-peak norm using an auxiliary system where \textit{impulsive disturbance} is replaced by an \textit{initial condition}. 
\begin{definition}\label{def:h2norm}
Given $\{\mcl T, \mcl A, \mcl B, \mcl C\}$, define the auxiliary PIE as 
%We say that System~\eqref{PIE2} has $H_2$ norm $\gamma$ if, for any $x_0 \in \R^{n_u}$ with $\norm{x_0}\le 1$ and $\mbf x, z$ which satisfy
\begin{align}\label{eqn:PIEaux}
\partial_t (\mcl T \mbf x(t))&= \mcl A \mbf x(t),\quad \mcl T \mbf{x}(0) =\mcl{B} v,\quad z(t)=\mcl C \mbf x(t),
\end{align}
where $v \in \R^{n_w}$. Then, the impulse-to-peak norm of system in Eq.~\eqref{PIE2} is given by
\[
 \norm{G(\mcl T, \mcl A,\mcl B,\mcl C)}_{ip}:=\sup_{\substack{z,\mbf x\, \text{satisfy~\eqref{eqn:PIEaux}}\\ \norm{v}=1}} \norm{z}_{L_\infty}.
\]
%we have that $\norm{z}_{L_2}\le \gamma$.
\end{definition}
\noindent In simpler words, impulse response of Eq. \eqref{PIE2} is equal to initial condition response of Eq. \eqref{eqn:PIEaux}. This equivalence can be established as follows: Let $w_n(t)\in L_1[\R_+]$ with $w_n(t)\to v\delta(t)$ (converges in a distributional sense). Assuming well-posedness (and by continuity of the input to state map \cite[Lemma 3.1.5]{curtain2012introduction}), we conclude that the sequence of weak solutions of \eqref{PIE2} converges to the weak solution of \eqref{eqn:PIEaux}.

Similar to Def. \ref{def:h2norm}, we can define an auxiliary PIE for the dual of Eq. \eqref{PIE2} as 
\begin{align}\label{eqn:PIEdaux}
\partial_t(\mcl T^* \bar{\mbf x})(t)&= \mcl A^* \bar{\mbf x}(t),\; \mcl T^* \bar{\mbf x}(0) =\mcl{C}^* \bar v,\;\bar z(t)=\mcl B^* \bar{\mbf x}(t),
\end{align}
 and corresponding I2P-norm of $G(\mcl T^*,\mcl A^*,\mcl C^*,\mcl B^*)$ as
\[
 \norm{G(\mcl T^*, \mcl A^*,\mcl C^*,\mcl B^*)}_{ip}:=\sup_{\substack{\bar{z},\bar{\mbf x}\, \text{satisfy~\eqref{eqn:PIEdaux}}\\ \norm{\bar v}=1}} \norm{\bar z}_{L_\infty}.
\]

\subsection{Equivalence in Impulse-to-peak Norm}
Since impulse response of a PIE and initial condition response of its auxiliary system have equivalent I2P-norms, we must now establish that initial condition responses of Eq. \eqref{eqn:PIEaux} and Eq. \eqref{eqn:PIEdaux} are equivalent---implying impulse response of $G(\mcl T,\mcl A,\mcl B,\mcl C)$ and $G(\mcl T^*,\mcl A^*,\mcl C^*,\mcl B^*)$ are equivalent.

%Having defined the impulse-to-peak norm, the goal now is to establish the duality, i.e., the equivalence of impulse-to-peak norms of the primal and dual systems.
\begin{theorem}\label{thm:dualityI2P}
    Given $\mcl{T,A}, \mcl B, \mcl C \in \Pi_4$, if the PIE Eq. \eqref{eqn:PIEaux} and its dual have well-posed solutions, then
    \[
    \norm{G(\mcl T, \mcl A,\mcl B,\mcl C)}_{ip}=\norm{G(\mcl T^*, \mcl A^*,\mcl C^*,\mcl B^*)}_{ip}.
    \]
\end{theorem}
\begin{proof}
    Suppose $\{\mbf x, z\}$ satisfy the PIE, Eq.~\eqref{eqn:PIEaux}, and $\{\bar{\mbf x}, \bar z\}$ satisfy the dual PIE, Eq.~\eqref{eqn:PIEdaux}. As in~\cite{shivakumar_dual_2024}, we can apply integration by parts to get
    \begin{align*}
       I_1:=& \int_0^t \ip{\bar{\mbf x}(t-s)}{\partial_s(\mcl T \mbf x(s))}_{\R L_2}ds
       \\&= \ip{\bar{\mbf x}(0)}{\mcl T \mbf x(t)}_{\R L_2}-\ip{\bar{\mbf x}(t)}{\mcl T \mbf x(0)}_{\R L_2}\\& \underbrace{-\int_0^t \ip{\partial_s\bar{\mbf x}(t-s)}{\mcl T \mbf x(s)}_{\R L_2} ds}_{:=I_2}.
    \end{align*}
Moreover, from the dynamics Eq. \eqref{eqn:PIEaux} and Eq. \eqref{eqn:PIEdaux}, we have
\begin{align*}
    I_1&=\int_0^t \ip{\bar{\mbf x}(t-s)}{\mcl A \mbf x(s)}ds 
    = \int_0^t \ip{\mcl A^*\bar{\mbf x}(t-s)}{\mbf x(s)}d\theta\\
    &= -\int_0^t \ip{\partial_s(\mcl T^*\bar{\mbf x})(t-s)}{\mbf x(s)}d\theta = I_2,
    %&=\int_0^t \ip{\partial_\theta(\mcl T^*\bar{\mbf x})(\theta)}{\mbf x(t-\theta)}_{\R L_2}d\theta\\
    %&=\int_0^t \ip{\partial_\theta\bar{\mbf x}(\theta)}{\mcl T\mbf x(t-\theta)}_{\R L_2}d\theta\\& = -\int_0^t \ip{\partial_s\bar{\mbf x}(t-s)}{\mcl T \mbf x(s)}_{\R L_2} ds= I_2.
\end{align*}
if $\bar{\mbf x}$ is differentiable.
% and
% \begin{align*}
%     I_2=\int_0^t \ip{\partial_{\theta}\bar{\mbf x}(\theta)}{\mcl T \mbf x(t-\theta)}_{\R L_2} d\theta
%     \\=
%     \int_0^t \ip{\mcl T^*\partial_{\theta}\bar{\mbf x}(\theta)}{\mbf x(t-\theta)}_{\R L_2} d\theta=I_1.
% \end{align*}
Therefore, $\ip{\mcl T^*\bar{\mbf x}(0)}{\mbf x(t)}_{\R L_2}=\ip{\bar{\mbf x}(t)}{\mcl T\mbf x(0)}_{\R L_2}$. Next, substituting the initial conditions $\mcl T^*\bar{\mbf x}(0) = \mcl C^* \bar v$ and $\mcl T \mbf x(0)=\mcl B v$, we get
\[
    \ip{\mcl C^* \bar v}{ \mbf x(t)}_{\R L_2}=\ip{\bar{\mbf x}(t)}{\mcl B v}_{\R L_2}.
\]
This implies, for any $t\ge 0$,
\begin{align*}
\ip{\bar v}{ z(t)}_{\R}&= \ip{\bar v}{ \mcl C \mbf x(t)}_{\R}=\ip{\mcl C^*\bar v}{ \mbf x(t)}_{\R L_2}\\
&=\ip{\bar{\mbf x}(t)}{\mcl Bv}_{\R L_2}=\ip{\mcl B^*\bar{\mbf x}(t)}{v}_{\R} =\ip{\bar{z}(t)}{v}_{\R}.
\end{align*}
For clarity, let us make explicit the dependence of $z$ and $\bar z$ on $v$ and $\bar v$, respectively, to rewrite the above identity as $\ip{\bar v}{z(t;v)}_{\R} = \ip{v}{\bar z(t;\bar v)}_{\R}$.

By definition of the vector norm, we have
\begin{align*}
    \sup_{\norm{\bar{v}}=1}\norm{\bar{z}(t;\bar{v})}_{\mathbb{R}} &= \sup_{\norm{\bar v}_2=1,\norm{v}=1}\!\!\!\!\ip{v}{\bar z(t;\bar v)}_{\R}\\
    & = \sup_{\norm{\bar v}_2=1,\norm{v}=1}\!\!\!\!\ip{\bar v}{z(t;v)}_{\R}\!\!=\!\! \sup_{\norm{v}=1}\norm{z(t;v)}_{\mathbb{R}}.
\end{align*}
Since this holds for any $t\in \R_+$, we have 
\begin{align*}
&\sup_{\norm{v}=1} \norm{z}_{\infty} = \sup_{\norm{v}=1}\left(\sup_{t\ge 0} \norm{z(t;v)}_{\mathbb{R}}\right)\\
&\quad =\sup_{t\ge 0} \left(\sup_{\norm{v}=1} \norm{z(t;v)}_{\mathbb{R}}\right)=\sup_{t\ge 0} \left(\sup_{\norm{\bar v}=1} \norm{\bar z(t;\bar v)}_{\mathbb{R}}\right)\\
&\quad =\sup_{\norm{\bar v}=1} \left(\sup_{t\ge 0} \norm{\bar z(t;\bar v)}_{\mathbb{R}}\right)= \sup_{\norm{\bar v}=1} \norm{\bar z}_{\infty}. 
\end{align*}
% Then,
% \[
% \ip{\bar v}{z(T;v)}_{\R}\ip{\bar v}{z(T;v)}_{\R} = \ip{v}{\bar z(T;\bar v)}_{\R}\ip{v}{\bar z(T;\bar v)}_{\R}.
% \]
% Taking supremum on both side, we get
% \begin{align*}
% &\sup_{\norm{\bar v}_2=1,\norm{v}=1}\ip{\bar v}{z(T;v)}_{\R}\ip{\bar v}{z(T;v)}_{\R} = \\&\sup_{\norm{\bar v}_2=1,\norm{v}=1}\ip{v}{\bar z(T;\bar v)}_{\R}\ip{v}{\bar z(T;\bar v)}_{\R}
% \end{align*}
% In other words, by property 
% \begin{align*}
% \norm{x} = \sup_{\norm{y}=1}\ip{y}{x}
% \end{align*}
% We can say 
% \begin{align*}
% \sup_{\norm{\bar{v}}=1}\norm{\bar{z}(t;\bar{v})}_{\mathbb{R}} = \sup_{\norm{v}=1}\norm{z(t;v)}_{\mathbb{R}}
% \end{align*}
% Now let us define the dual generalized $H_2$ norm:
% \begin{align*}
% \norm{\mathcal{G}_{dual}}^2_{gen} = \sup_{\norm{\bar{v}} = 1}(\sup_{t \geq 0}\norm{\bar{z}(t;\bar{v})}^2_{\mathbb{R}}) = \gamma^2
% \end{align*}
% We can swap sup
% \begin{align*}
% \norm{\mathcal{G}_{dual}}^2_{gen} = \sup_{t \geq 0}(\sup_{\norm{\bar{v}} = 1}\norm{z(t;\bar{v})}^2_{\mathbb{R}})
% \end{align*}
% By above property
% \begin{align*}
% \norm{\mathcal{G}_{dual}}^2_{gen} =\sup_{t \geq 0}(\sup_{\norm{v} = 1}\norm{z(t;v)}^2_{\mathbb{R}})
% \end{align*}
% Then, we can swap sup back
% \begin{align*}
% \norm{\mathcal{G}_{dual}}^2_{gen} =\sup_{\norm{v} = 1}(\sup_{t \geq 0}\norm{z(t;v)}^2_{\mathbb{R}})
% \end{align*}

% \begin{align*}
% \norm{\mathcal{G}_{dual}}^2_{gen} =\sup_{\norm{v} = 1}(\sup_{t \geq 0}\norm{z(t;v)}^2_{\mathbb{R}})
% \end{align*}
% Rewrite
% \begin{align*}
% \norm{\mathcal{G}_{dual}}^2_{gen} =\sup_{\norm{v} = 1}(\norm{z}_{L_\infty}^2)
% \end{align*}
% This is equal to both $\gamma^2$ and the primal $H_2$ generalized norm.
% Thus 
% \begin{align*}
% \norm{\mathcal{G}_{dual}}^2_{gen} =\gamma^2 = \norm{\mathcal{G}}^2_{gen}
% \end{align*}
\end{proof}


\section{Linear PI Inequality Formulations}

Having established the duality result for the impulse-to-peak norm of PIEs, we will next present a convex optimization based characterization of the norm using Lyapunov functions. 
\begin{theorem}\label{thm:Normcoercive}
    Given 4-PI operators ${\mathcal{T},\mathcal{A},\mathcal{B},\mathcal{C}}$, let $\{\mbf x, z\}$ satisfy a PIE of the form
    Eq. \eqref{eqn:PIEaux}. Suppose there exists a $\gamma>0$ and a 4-PI operator $\mcl Q \succeq 0$ such that one of the following two conditions hold:
    \begin{enumerate}[leftmargin=*]
        \item[(i)] $\mathcal{B}^*\mcl Q\mathcal{B}\preceq I$, $\mathcal{A}^*\mcl Q\mathcal{T}+\mathcal{T}^*\mcl Q\mathcal{A} \preceq 0$, $\frac{1}{\gamma^2}\mathcal{C}^*\mathcal{C} \preceq \mathcal{T}^*\mcl Q\mathcal{T}$.
        \item[(ii)] $\mathcal{C}\mcl Q\mathcal{C}^*\preceq I$, $\mathcal{A}\mcl Q\mathcal{T}^*+\mathcal{T}\mcl Q\mathcal{A}^* \preceq 0$, $\frac{1}{\gamma^2}\mathcal{B}\mathcal{B}^* \preceq \mathcal{T}\mcl Q\mathcal{T}^*$.
    \end{enumerate} 
    Then, $\norm{G(\mcl T, \mcl A,\mcl B,\mcl C)}_{ip} \le \gamma$.
\end{theorem}

\begin{proof}[Proof of (i)]
     Let $\{\mbf x, z\}$, satisfy Eq. \eqref{eqn:PIEaux} with initial conditions $\mcl T\mbf x(0) = \mcl Bv$. Define a storage function $V(\mbf x) = \langle\mathcal{T}\mbf x, \mcl Q\mathcal{T}\mbf x\rangle$ where $\mcl Q\succeq 0$ satisfies constraints in \textit{(i)}. Then, $V(0) = \ip{\mcl B v}{\mcl Q\mcl Bv} = \ip{v}{\mcl B^*\mcl Q\mcl B v} \leq \norm{v}^2$. Differentiating $V$ along the solutions of Eq. \eqref{eqn:PIEaux}, we get 
\begin{align*}
&\dot{V} = \langle \partial_t \mathcal{T}\mbf x, \mcl Q \mathcal{T}\mbf x\rangle+\langle \mathcal{T}\mbf x, \mcl Q \partial_t \mathcal{T}\mbf x\rangle\\
&= \langle \mathcal{A}\mbf x, \mcl Q \mathcal{T}\mbf x\rangle\!+\!\langle \mathcal{T}\mbf x, \mcl Q \mathcal{A}\mbf x\rangle \!\! =\!\! \langle \mbf x, (\mathcal{A}^* \mcl Q \mathcal{T}\!+\!\mathcal{T}^*\mcl Q\mathcal{A})\mbf x\rangle\!\le\! 0.    
\end{align*}
Lastly, we have that 
\begin{align*}
\frac{1}{\gamma^2}\norm{z(t)}^2&=\frac{1}{\gamma^2}\ip{\mcl C\mbf x(t)}{\mcl C\mbf x(t)}=\ip{\mbf x(t)}{\frac{1}{\gamma^2}\mcl C^*\mcl C\mbf x(t)}\\
&\le\ip{\mbf x}{\mcl T^*\mcl Q\mcl T\mbf x} = V(\mbf x(t)).
\end{align*}
Thus, $\lVert z(t) \rVert^2 \leq \gamma^2 V(t) \leq \gamma^2 V(0) \leq \gamma^2 \lVert v\rVert^2$ for all $t>0$ and hence, 
\[
\norm{G(\mcl T,\mcl A,\mcl B,\mcl C)}_{ip}=\sup_{\norm{v}=1}\norm{z}_\infty=\sup_{\lVert v\rVert = 1} \sup_{t\geq 0} \lVert z(t) \rVert \leq \gamma.\]
Alternatively, let $\mcl Q$ satisfy the constraints in \textit{(ii)}. Then, following the steps above one can show that $\norm{G(\mcl T^*,\mcl A^*,\mcl C^*,\mcl B^*)}_{ip}\le \gamma$ where $\bar{\mbf x}, \bar z$ satisfy Eq. \eqref{eqn:PIEdaux} with initial conditions $\mcl T^*\bar{\mbf x}(0)=\mcl C^*\bar v$. From Thm. \ref{thm:dualityI2P}, the primal also satisfies the same bound.
\end{proof}

%\textcolor{blue}{While the feasibility problem in Thm. \ref{thm:Normcoercive} provides sufficient conditions to find an upper bound on the impulse-to-peak norm, these can be conservative. To alleviate this conservatism, we use a broader class of Lyapunov functions to derive the following LPI for upper bounding the norm. }

To alleviate conservatism in the sufficient conditions given in Thm \ref{thm:Normcoercive}, we use a broader class of Lyapunov functions to derive the following LPI for upper bounding the norm. 

\begin{theorem}\label{thm:Normnoncoercive}
    Given 4-PI operators ${\mathcal{T}, \mathcal{A},\mathcal{B},\mathcal{C}}$, let $\{\mbf x, z\}$ satisfy a PIE of the form Eq. (\ref{eqn:PIEaux}).
    Suppose there exists a $\gamma > 0$ and a 4-PI operator $\mcl Q$ such that $\mcl T^*\mcl Q =\mcl Q^* \mcl T \succeq 0$ and one of the following conditions hold:
    \begin{enumerate}
        \item[(i)] $\mcl A^*\mcl Q+\mcl Q^*\mcl A \!\preceq\! 0$, $ \begin{bmatrix}
            \gamma^2I & \mcl C\\ \mcl C^* & \mcl Q^* \mcl  T 
        \end{bmatrix}\!\succeq\! 0$, $ \begin{bmatrix}
            \mcl T^*\mcl Q & \mcl Q^* \mcl B\\ \mcl B^* \mcl Q & I
        \end{bmatrix}\!\succeq\! 0$ 
        \item[(ii)] $\mcl A\mcl Q+\mcl Q^*\mcl A^* \!\preceq\! 0$, $ \begin{bmatrix}
            \gamma^2I \!\!& \mcl B^*\\ \mcl B \!\!& \mcl Q^* \mcl  T^*
        \end{bmatrix}\!\succeq\! 0$, $ \begin{bmatrix}
            \mcl T\mcl Q & \mcl Q^* \mcl C^*\\ \mcl C \mcl Q & I
        \end{bmatrix}\!\succeq\! 0$ 
    \end{enumerate} 
    Then, $\norm{G(\mcl T, \mcl A,\mcl B,\mcl C)}_{ip} \le \gamma$.
\end{theorem}

\begin{proof}
We give a proof outline, since it is similar to the proof of Thm. \ref{thm:Normcoercive}.
Let $\{\mbf x, z\}$ satisfy the PIE Eq. \eqref{eqn:PIEaux} with initial conditions $\mcl T\mbf x(0)=\mcl Bv$. Next, define a storage function as $V(\mbf x) = \ip{\mcl T\mbf x}{\mcl Q\mbf x}$. Then, if $\mcl Q$ satisfies the constraints in \textit{(i)}, we have
$\dot{V}(\mbf x(t))\le 0$ and $\norm{z(t)}^2\le \gamma^2V(t)$ for all $t>0$.
% \begin{align*}
%     V(\varphi) = \ip{\varphi}{\mcl T^* \mcl Q \varphi} = \ip{\varphi}{\mcl Q^* \mcl T \varphi}\\
%     \dot{V}(\varphi) = \ip{\varphi}{\mcl Q^* \partial_t \mcl T \varphi} + \ip{\partial_t \varphi}{\mcl Q^* \mcl T \varphi}\\
%     = \ip{\varphi}{\mcl Q^* \partial_t \mcl T \varphi} + \ip{\mcl Q^*\mcl T\partial_t \varphi}{\varphi}\\
%     =\ip{\varphi}{\mcl Q^* \mcl A \varphi} + \ip{\mcl Q^*\mcl A \varphi}{\varphi}\\ 
%     = \ip{\varphi}{\mcl Q^* \mcl A \varphi}+\ip{\varphi}{\mcl A^* \mcl Q \varphi}\\
%     = \ip{\varphi}{(\mcl Q^* \mcl A+\mcl A^*\mcl Q)\varphi}
% \end{align*}
% Using assumption 1, the above storage function, with negative semi-definite time derivative, is non-increasing in time. 
% \begin{comment}
       
    
%         \item Take the Lyapunov function $V(\varphi) = \ip{\varphi}{\mcl R \varphi} = \ip{\varphi}{\mcl Q^* \mcl T \varphi} = \ip{\mcl Q \varphi}{\mcl T \varphi}$. We can rewrite the first part of the inner product such that the total expression gives $\ip {\mcl Q \mcl T^{-1} \mcl T \varphi}{\mcl T \varphi}= \ip{\mcl Q \mcl T^{-1} x}{x}$, where in this proof $x$ represents the PDE state. Taking the time derivative of the storage function, $D_tV = \ip{\mcl Q \mcl T^{-1} \partial_t x}{x}+ \ip{\mcl Q \mcl T^{-1} x}{\partial_tx} = \ip{\mcl Q \mcl T^{-1} \mcl A \varphi}{x}+\ip{\mcl Q \mcl T ^{-1}x}{\mcl A \varphi}$. The second inner product term can utilize $x = \mcl T \varphi$ to map $\mcl T ^{-1} x$ back to $\varphi$. Thus the expression becomes $D_tV = \ip{\varphi}{\mcl A \mcl T^{-*}\mcl Q ^*x}+ \ip{\varphi}{\mcl Q ^* \mcl A \varphi}$. Because $\mcl Q^*\mcl T = \mcl T^* \mcl Q \implies \mcl T^{-*}\mcl Q^*\mcl T = \mcl T^{-*}\mcl T^*Q = Q.$ Thus, $D_tV = \ip{\varphi}{\mcl A^* \mcl T^{-*}\mcl Q^* \mcl T \varphi}+\ip{\varphi}{\mcl Q^* \mcl A \mcl T} = \ip{\varphi}{\mcl A^* \mcl Q  \varphi}+\ip{\varphi}{\mcl Q^* \mcl A \mcl T} = \ip{\varphi}{(\mcl A^* \mcl Q + \mcl Q^*\mcl A)\varphi}$. Thus, if $\mcl A^* \mcl Q + \mcl Q^*\mcl A \preceq 0$, the overall storage function time derivative will be nonpositive.
% \end{comment}

        % Next, $\norm{z(t)}^2 = \ip{z(t)}{z(t)} = \ip{\mcl C \varphi}{\mcl C \varphi} = \ip{\varphi}{\mcl C^*\mcl C\varphi}.$ If the second condition is met, by Schur Complement, $\mcl C^* \mcl C \preceq \mcl Q^* \mcl T$, then $\ip{\varphi}{\mcl C^* \mcl C\varphi}\preceq \gamma^2 \ip{\varphi}{\mcl Q^* \mcl T\varphi}= \gamma^2 \ip{\mcl Q \varphi }{\mcl T \varphi} = \gamma^2 V(t)$
        % Finally, we show that the storage function at time 0 is bounded by the initial storage function. $V = \ip{\mcl Q\varphi}{\mcl T\varphi}$. Then using $\mcl T \varphi = \mcl B v$, 
        Lastly, for any $v\in \R^{n_w}$ and $\mcl T\mbf x(0)=\mcl Bv$, we have
        \begin{align*}
            0\le &\ip{\bmat{\mbf x(0)\\-v}}{ \bmat{\mcl T^*\mcl Q&\mcl Q^*\mcl B\\\mcl B^*\mcl Q&I} \bmat{\mbf x(0)\\-v}} \\
            &= \ip{\mbf x(0)}{\mcl T^*\mcl Q\mbf x(0)} - 2\ip{\mbf x(0)}{\mcl Q^*\mcl Bv} + \norm{v}^2 \\
            &=\ip{\mbf x(0)}{\mcl Q^*\mcl T\mbf x(0)} - 2\ip{\mbf x(0)}{\mcl Q^*\mcl T\mbf x(0)} + \norm{v}^2 \\
            &= \norm{v}^2-V(\mbf x(0)).
        \end{align*}
        Since $\norm{z(t)}^2\le \gamma^2V(\mbf x(t))\le \gamma^2V(\mbf x(0))\le \gamma^2\norm{v}^2$, we have $\norm{G(\mcl T, \mcl A,\mcl B,\mcl C)}_{ip} \le \gamma$.
        Likewise, satisfaction of \textit{(ii)} implies the norm bound on the dual and the primal PIE response.
       % \begin{align*}
        % &V(\mbf x(0)) \!= \!\ip{\mcl Q \mbf x(0)}{\mcl B v} \!= \!\ip{\mbf x(0)}{\mcl Q^* \mcl B v} \!= \!\ip{\mcl T^{-1}\mcl T\mbf x(0)}{\mcl Q^* \mcl B v} \\
        % &= \ip{\mcl T ^{-1}\mcl B v}{\mcl Q ^ * \mcl Bv}= \ip{v}{\mcl B ^* \mcl T ^{-*} \mcl Q^* \mcl B v}
        % \end{align*}
        % Using assumption from the theorem,$ \ip{v}{\mcl B ^* \mcl T ^{-*} \mcl Q^* \mcl B v} \preceq \ip{v}{Iv} = \ip{v}{v} = \norm{v(t)}^2$. Thus, $V(0) \leq \norm{v^2}$. 
        % The final argument chains together similarly to the previous theorem. The last condition tells you that the initial storage function is bounded, the first condition tells you that the storage function decreases, and the second condition tells you the output is also bounded by the storage function scaled. We can chain everything together as $\norm{z(t)}^2\leq \gamma^2V(t) \leq \gamma^2 V(0)\leq \gamma^2 \norm{v}^2 $. Taking initial condition with norm 1, then $\norm{z(t)}^2 \leq \gamma^2.$ Finally $\sup_{\norm{v}= 1}\sup_{t\geq 0}\norm{z(t)}^2\leq \gamma^2$. 
\end{proof}

% Note that the equivalence of the condition $\mcl B^*\mcl T^{-*}\mcl Q^* \mcl B \preceq I$ with $\begin{bmatrix}
%     \mcl R & \mcl Q^* \mcl B\\ \mcl B^* \mcl Q &I
% \end{bmatrix} \succeq 0 $ can be explained by writing $\mcl T^{-*} \mcl Q^*$ as $\mcl Q \mcl Q^{-1}\mcl T^{-*} \mcl Q^*  = Q(\mcl T^* \mcl Q)^{-1} \mcl Q^*$ making the whole expression $\mcl B^* Q(\mcl T^* \mcl Q)^{-1} \mcl Q^*\mcl B  -I \preceq 0$. Applying Schur Complement, and taking $\mcl T^* \mcl Q$ as $\mcl R$ yields the block condition.

\subsection{Optimal State-feedback Controller Synthesis}
Using the duality results and a Lyapunov function-based characterization of the norm, we can consider the I2P optimal control problem. Specifically, given a PIE of the form,
%Now that we have established duality results and a Lyapunov function-based characterization of the norm, we can consider the impulse-to-peak optimal control problem. Specifically, given a PIE of the form,
\begin{align}
    \partial_t(\mcl T\mbf x)(t) &= \mcl A\mbf x(t) + \mcl B w(t) + \mcl B_2u(t), \quad \mcl T\mbf x(0) = 0,\notag\\
    z(t) &= \mcl C\mbf x(t) + \mcl D u(t), \label{eq:PIEcontrol}
\end{align}
the goal is to design a state-feedback $u(t) = \mcl K\mbf x(t)$ such that impulse-input to peak-output norm is optimal. %i.e., find $\mcl K$ minimizing $\gamma$ such that $\sup_{\norm{v}=1} \norm{z}_\infty\le \gamma$, where $w(t) = \delta(t)v$. 

Substituting the control policy, the closed-loop system takes the form $G(\mcl T, \mcl A+\mcl B_2\mcl K,\mcl B, \mcl C+\mcl D\mcl K)$ where $G$ is defined in Definition \ref{def:ip-normsystem}. Thus, the control task can be stated as: Given a closed-loop system $G$, find $\mcl K$ minimizing $\gamma$ such that $\sup_{\norm{v}=1} \norm{z}_\infty\le \gamma$.  

Note the auxiliary system admits the form
\begin{align}\label{eqn:PIEauxwinput}
\partial_t(\mcl T\mbf x)(t) &= (\mcl A+\mcl B_2\mcl K)\mbf x(t), \quad\mcl T\mbf x(0) = \mcl Bv, \notag\\
z(t) &= (\mcl C+\mcl D\mcl K)\mbf x(t).
\end{align}
% whose auxiliary system has the form
% \begin{align}
%     \partial_t (\mcl T \mbf x(t))&= (\mcl A +\mcl B_2\mcl K)\mbf x(t),\quad \mcl T \mbf{x}(0) =\mcl{B} v,\\
%     z(t)&=(\mcl C +\mcl D\mcl K)\mbf x(t),\notag 
% \end{align}

% Under a state feedback policy $\mbf u(t) = \mcl K \mbf x(t)$, the system defined in equation \ref{eqn:PIEauxwinput} becomes 
% \begin{align}\label{eqn:PIECL}
%     \partial_t (\mcl T \mbf x(t))&= (\mcl A + \mcl B_u \mcl K) \mbf x(t),\quad \mcl T \mbf{x}(0) =\mcl{B} v,\\
%     z(t)&=\mcl C \mbf x(t)
% \end{align}.

% We call equation \ref{eqn:PIECL} the closed loop system auxillary PIE. This leads to the final definition.
% \begin{definition}
%     The closed loop system dual auxillary PIE is given by 
%     \begin{align}\label{eqn:PIECLdual}
%     \partial_t (\mcl T^* \bar{\mbf x}(t))&= (\mcl A^* + \mcl K^* \mcl B_u^* ) \bar{\mbf x}(t),\quad \mcl T^* \mbf{\bar{x}}(0) =\mcl{C}^* \bar{v},\\
%     \bar{z}(t)&=\mcl B^* \bar{\mbf x}(t)
% \end{align}
% \end{definition}

If we apply Thm. \ref{thm:Normcoercive} to the auxiliary system defined above, the constraints in \textit{(i)} are bilinear in $\mcl Q$ and $\mcl K$ and intractable. However, contingent on the change of variable $\mcl Z=\mcl K \mcl Q$, the constraints in \textit{(ii)} are linear and convex. Since \textit{(i)} and \textit{(ii)} are equivalent, we can solve \textit{(ii)} to search for a controller $\mcl K = \mcl Z\mcl Q^{-1}$ that achieves optimal $\gamma$ for the closed-loop system.


% As our LPI developed in Theorem 2 is designed for system of this form, we can develop an LPI for the state feedback controller design.

% Consider dual auxillary closed loop system with input defined in equation \ref{eqn:PIECLdual}.
\begin{corollary}\label{cor:control}
    Given 4-PI operators $\{\mathcal{T}, \mathcal{A},\mathcal{B}, \mcl B_2,\mathcal{C}, \mcl D\}$, let $\{\mbf x, z\}$ satisfy a PIE of the form Eq. (\ref{eqn:PIEauxwinput}).
    Suppose there exists a $\gamma > 0$
    such that one of the following sets of conditions hold:
    \begin{enumerate}[leftmargin=6mm]
        \item[(i)] $\mcl Q\succ 0$, $\mcl A\mcl Q \mcl T^*+\mcl B_2 \mcl Z \mcl T^* + \mcl T \mcl Q\mcl A^* + \mcl T \mcl Z^* \mcl B_2^* \preceq 0,\vspace{1mm}\\  \begin{bmatrix}
            \gamma^2I&  \mcl B^*\\ \mcl B & \mcl T\mcl Q\mcl T^*
        \end{bmatrix}\!\succeq\! 0,\;\bmat{I & (\mcl C\mcl Q+\mcl D\mcl Z)\\(\mcl C\mcl Q+\mcl D\mcl Z)^*& \mcl Q}\!\succeq\! 0$. \vspace{3mm}
       
        \item[(ii)] $\mcl Q^*\mcl T^*=\mcl T\mcl Q$, $\mcl A \mcl Q +\mcl B_2 \mcl Z +\mcl Q^* \mcl A^* + \mcl Z^* \mcl B_2^* \preceq 0,\vspace{1mm}\\ \begin{bmatrix}
            \gamma^2I & \mcl B^*\\ \mcl B & \mcl Q^* \mcl  T^* 
        \end{bmatrix}\!\succeq\! 0, \;\begin{bmatrix}
            I & (\mcl C\mcl Q+\mcl D\mcl Z)\\  (\mcl C\mcl Q+\mcl D\mcl Z)^*& \mcl Q^*\mcl T^*
        \end{bmatrix}\!\succeq\! 0$. 
    \end{enumerate}
    Then, under the feedback control $u(t)=\mcl K\mbf x(t)$ where $\mcl K= \mcl Z\mcl Q^{-1}$, we have $\sup_{\norm{v}=1} \norm{z}_\infty \le \gamma$.
\end{corollary}
\begin{proof}
The proof follows the exact same steps as proofs of \textit{(ii)} in Thms. \ref{thm:Normcoercive} and \ref{thm:Normnoncoercive}, with the additional change of variable $\mcl K\mcl Q = \mcl Z$. If there exists a $\gamma>0, \mcl Q$ satisfying \textit{(i)} or \textit{(ii)} above, and the solution $\{\mbf x, z\}$ satisfies the closed-loop PIE under the feedback $u(t)=\mcl K\mbf x(t)$, then $\gamma, \mcl Q$ satisfy the conditions in \textit{(ii)} of Thm. \ref{thm:Normcoercive} or \ref{thm:Normnoncoercive}, respectively, for the closed-loop PIE. Thus, the closed-loop PIE Eq. \eqref{eqn:PIEauxwinput} satisfies the norm bound.
% The proof for the first conditions relies on utilizing the LPIs in Theorem 2.

% Specifically, condition 1 becomes:
% \begin{align}
%     (\mcl A + \mcl B_u\mcl K)\mcl Q \mcl T^*+ \mcl T\mcl Q(\mcl A^* +\mcl K^*\mcl B_u^*) \preceq 0\\
%     \iff (\mcl A + \mcl B_u \mcl K)\mcl Q\mcl T^* +\mcl T\mcl Q\mcl A^*+\mcl T\mcl Q\mcl K^* \mcl B_u^* \preceq 0\\
%     \iff \mcl A\mcl Q \mcl T^*+\mcl B_u \mcl K \mcl Q \mcl T^* + \mcl T \mcl Q\mcl A^* + \mcl T \mcl Q \mcl K ^* \mcl B_u^* \preceq 0 \label{presubstitute}
% \end{align}

% We then make the substitution $\mcl K\mcl Q = \mcl Z$.

% This simplifies the above expression in eq. \ref{presubstitute} to:
% \begin{align}
%     \mcl A\mcl Q \mcl T^*+\mcl B_u \mcl Z \mcl T^* + \mcl T \mcl Q\mcl A^* + \mcl T \mcl Z^* \mcl B_u^* \preceq 0 
% \end{align}

% This condition is now convex in $\mcl Z$ and $\mcl Q$.

% The proof for the second set of conditions relies on utilizing the LPIs in Theorem 3.

% Specifically, condition 1 becomes:
% \begin{align}
%     ((\mcl A^* +\mcl K^*\mcl B_u ^*)^*\mcl Q + \mcl Q^*(\mcl A^*+ \mcl K^*\mcl B_u ^*)) \preceq 0 \implies \\ 
% (\mcl A\mcl Q +\mcl B_u\mcl K\mcl Q  + \mcl Q^*\mcl A^*+ \mcl Q^*\mcl K^*\mcl B_u^* ) \preceq 0 
% \end{align}

% We then make the substitution $\mcl K\mcl Q = \mcl Z$.

% This simplifies the above expression in eq. \ref{presubstitute} to:
% \begin{align}
%     \mcl A \mcl Q +\mcl B_u \mcl Z +\mcl Q^* \mcl A^* + \mcl Z^* \mcl B_u^* \preceq 0 
% \end{align}

% This condition is now convex in $\mcl Z$ and $\mcl Q$.
\end{proof}


\section{Numerical Results}

In this section, we present numerical examples validating the above theoretical results. First, we present two PDE examples for which I2P norm bounds are computed using the LPIs described in Theorem \ref{thm:Normnoncoercive}. For the state feedback control design, we use the conditions in \textit{(i)} of Cor. \ref{cor:control}. All the steps required to solve the optimization problems in Thms. \ref{thm:Normnoncoercive} and Cor. \ref{cor:control}, are implemented using the PIETOOLS toolbox.%such as setting up, solving, and controller gain reconstruction via operator inverse, 


% The underlying optimization problems are solved using an interior point solver in the case of the open loop norm bounding, and a standard bisection method is used for the controller design and closed loop norm bounding.
\subsection{Impulse-to-peak Norm Upper-bounding}
\begin{example}

Consider the following transport PDE:
\begin{align*}
  \partial_t x(t,s) &= \partial_sx(t,s)+(s-s^2)w(t),\; z(t) = \int_0^1 x(t,s)\ ds
\end{align*}
where $z$ is the regulated output and Dirichlet boundary condition $x(t,1) = 0$.
For this PDE, we use PIETOOLS to obtain the PIE parameters and solve the optimization problems presented in Thm. \ref{thm:Normnoncoercive} with $\gamma^2$ as the objective to be minimized. %The true bound is recovered after obtaining $\gamma_1$ as $\gamma = \sqrt{\gamma_1}$.
The upper bounds on I2P norm obtained from the primal and dual LPI solutions is $\gamma = 0.1684$ and $\gamma = 0.1667$ respectively, whereas the theoretical bound is $1/6$.
\end{example}

\begin{example}
Next, we consider the heat equation
\begin{align*}
\partial_t x(t,s) &= \partial^2_sx(t,s)+sw(t),\quad
z(t) = \int_0^1 x(t,s) \ ds,
\end{align*}
with boundary conditions $x(t,0) = \partial_s x(t,1) = 0$. We again solve the LPI problems in Thm. \ref{thm:Normnoncoercive} minimizing $\gamma^2$ and obtain $\gamma = 0.5000$, which matches the theoretical bound $1/2$. 
\end{example}

\subsection{Impulse-to-peak Norm Optimal State Feedback Control}

% \subsubsection{Class of Systems Where Controlled Norm will not Decrease}
% For naturally stable systems, it is possible that the initial condition of the system results in the peak norm having maximum value at $t = 0$. For example, for a transport equation, if the initial condition function has only positive values with one end fixed at 0, the integral of the PDE state across the domain (which can be taken as the regulated output) will always immediately decrease. 

% Specifically, take:
% \begin{align}
%     \partial_tx(t,s) = \partial_{s}x(t,s) + sw(t) +u(t)
% \end{align} 
% with output:
% \begin{align}
%     z(t) = \int_0^1 x(t,s)\ ds  
% \end{align}
% and boundary condition:
% $$x(t, 1) = 0$$

% Because of the zero boundary condition, if the spatial function is greater than zero for all s, the output can only reduce for any $t > 0$. The initial output is $\int_0^1 s\ ds = 1/2$, and thus the optimal $\gamma = 0.5$. Even with control, this optimal $\gamma$ value will remain the same.
% \subsubsection{Examples Where Control is Effective}
% This section presents two examples where control is effective. 
The first example considers I2P optimal control and stabilization of an unstable reaction-diffusion equation. In the second example, we consider a beam example and design a controller to reduce the peak output norm.

\begin{example}[Unstable Reaction-Diffusion Equation]
Consider an unstable reaction-diffusion equation:
\begin{align*}
    \partial_tx(t,s) &= 14 x(t,s)+ \partial_{s}^2x(t,s) + (s^2-2s)w(t)+u(t),\\
    z(t) &=  \int_0^1 2x(t,s)\ ds,\;\;\; x(t, 0) = \partial_s x(t,1) = 0.
\end{align*}
% It is straightforward to show using standard techniques that the above PDE is unstable for values of $\lambda > \pi^2$. 
Since this PDE is unconditionally unstable, the I2P gain is infinity. % for any non-trivial impulse input. 
To find an optimal stabilizing control, we solve the LPI presented in $\textit{(i)}$ of Cor. \ref{cor:control} for fixed $\gamma$ and apply bisection method to obtain the best achievable $\gamma$ and the corresponding controller. In this case, we can find a stabilizing controller that achieves $\gamma =  1.375$. The closed-loop system response with the obtained controller is shown in Fig. \ref{fig:reactevol}. Additionally, the regulated output is shown in Fig. \ref{fig:reactreg}.
\end{example}

\begin{figure}
    \centering
    \includegraphics[width=0.99\linewidth]{newestfigures/ConsistentFigs/ReactDiffuseSurf.pdf}
    \caption{Evolution of the PDE state for the reaction-diffusion equation under state-feedback controller}
    \label{fig:reactevol}
    \vspace{-15pt}
\end{figure}
\begin{figure}
    \centering
    \includegraphics[width=0.75\linewidth]{newestfigures/ConsistentFigs/ReactDiffuse.pdf}
    \caption{The regulated output, $\int_0^12 x(t,s) ds$, for the reaction-diffusion equation. The control stabilizes the system and satisfies peak output bounds.}
    \label{fig:reactreg}
\end{figure}

\begin{example}[Timoshenko Beam Equations]
Considering a beam with transverse displacement $w$ and cross-section rotation angle $\phi$, the Timoshenko beam equations model the evolution of these quantities with the following equations:
\begin{align}
\ddot{w} = \partial_s(w_s - \phi)+u+sw, \;\; \ddot{\phi} = \phi_{ss} +(w_s - \phi)  \nonumber 
\\ \phi(0) = w(0) = 0, \; \phi_s(0) = 0, \; w_s(1) - \phi(1) = 0. \nonumber
\end{align}

Then, we can choose $\mbf x = \begin{bmatrix}
    \dot{w},  w_s -\phi, \dot{\phi}, \phi_s
\end{bmatrix}$ and get:
\begin{align}
    \dot{\mbf x}(t,s) = \begin{bmatrix}
        0 & 0&0&0\\ 0&0&-1&0\\ 0 &1 &0&0\\ 0&0&0&0
    \end{bmatrix} \mathbf{x}(t,s) + \begin{bmatrix}
        0&1&0&0\\ 1&0&0&0\\ 0&0&0&1\\ 0&0&1&0
    \end{bmatrix} \partial_s \mbf{x}(t,s)\nonumber\\+\begin{bmatrix}
        1&0&0&0
    \end{bmatrix}^\top u + \begin{bmatrix}
        s&0&0&0
    \end{bmatrix}^\top w\nonumber,    
    %\\ \begin{bmatrix}
    %    1&0&0&0&0&0&0&0\\ 0&0&1&0&0&0&0&0\\ 0&0&0&0&0&0&0&1\\0&0&0&0&0&1&0&0
    %\end{bmatrix}
    \\\mathbf{x}_1(t,0) = \mathbf{x}_3(t,0) = \partial_s\mathbf{x}_4(t,0) = \partial_s\mathbf{x}_2(t,0) = 0 \nonumber
    %\begin{bmatrix}
    %    \mbf x(t,0)\\ \partial_s \mbf x(t,0)
    %\end{bmatrix} = 0 \nonumber
\end{align} with output  $z(t) = \int_0^1 \begin{bmatrix}
        1 & 0&0&0
    \end{bmatrix} \mathbf{x}(t,s) ds$. 

The closed-loop I2P-norm is 0.50, while the open-loop I2P-norm is 0.54, proving that controller suppresses peak-output norm.
\end{example}
\begin{figure}
    \centering
    \includegraphics[width=0.75\linewidth]{newestfigures/ConsistentFigs/BeamReg.pdf}
    \caption{The regulated output, $\int_0^1 \mathbf{x}_1(t,s)ds$, for the controlled ($\textcolor{red}{--}$) and uncontrolled ($\textcolor{blue}{-}$) Timoshenko beam equations.}
    \vspace{-14pt}
    \label{fig:placeholder}
\end{figure}

\begin{figure}
    \centering
    \vspace{-15pt}
    \includegraphics[width=0.99\linewidth]{newestfigures/ConsistentFigs/w.pdf}
    \caption{Evolution of ${w}$, for the Timoshenko Beam Equations under state-feedback control.}
    \label{fig:placeholder}
    \vspace{-15pt}
\end{figure}

\subsection{A Note on Computation}

Table \ref{tab:execution_time_comparison} provides the time per bisection iteration for the state-feedback control design problems we considered. It also displays the size of the underlying LMIs.

\begin{table}[h] 
\centering 
\begin{tabular}{lcccc} 
\textbf{\makecell{Model\\Type}} & \textbf{\makecell{Time \\(in CPU sec.)}} & \textbf{\makecell{\# of \\iterations}}& \textbf{\makecell{Decision \\Variables}} & \textbf{Constraints} \\ 
\midrule
Beam & 14.59 & 15 & 203,065 & 4,353 \\ 
\makecell{Reaction-\\diffusion} & 11.89 & 5 & 183,840& 1,054\\
\end{tabular} 
\caption{Computation Parameter and Time Comparison}
\vspace{-20pt}
\label{tab:execution_time_comparison} 
\end{table}


\section{Conclusion}
Despite its relevance in safe operations, a computational method to characterize I2P norm for a sufficiently general class of PDEs was not available. In this work, we leveraged the PIE representation to develop Lyapunov-based convex optimization formulation of the I2P norm bounding problem for arbitrary linear PDEs. We showed the bounds so obtained are provable in the sense the feasibility of LPIs provide a certificate for the I2P-norm bound. We also demonstrated there is no theoretical conservatism in the primal and dual I2P-norms. We solved the impulse-to-peak norm upper-bounding and norm-optimal state-feedback control synthesis problem for linear PDEs and validated the solution for various numerical examples, for which analytical bounds are known.





% Impulse-to-peak norm is frequently used to analyze the limits of safe operation for dynamical systems. However, a computational method to compute impulse-to-peak norm that applies to a sufficiently general class of PDEs was not available. This work solves this problem by leveraging the PIE framework for PDEs and formulating impulse-to-peak norm analysis and norm optimal state-feedback control problems as tractable convex optimization problems. Firstly, we showed that the PIE and dual PIE representation of a PDE have equivalent impulse-to-peak norms. Then, we proposed a Lyapunov function based computational test to find the smallest upper-bound on the impulse-to-peak norm. Coupling the computational test with the duality, we arrived at a convex optimization formulation of the impulse-to-peak optimal control design problem for PDEs represented as PIEs. We verified and validated the approach by applying the results for various numerical examples, for which analytical bounds on impulse-to-peak norm can be computed. Lastly, we numerically demonstrated that the controllers designed using this method can indeed stabilize and have provable performance guarantees.

% a proof for the equivalence of the impulse-to-peak norm for a PIE and its dual. Additionally, conditions are found for numerically bounding the impulse-to-peak norm from above. Using the duality results, the bounding conditions are able to be reformulated for controller design, that can stabilize unstable systems, or improve the impulse-to-peak upper bound for systems that are naturally stable open loop. This was demonstrated in numerical examples. Future extensions of this work will demonstrate the applicability of these methods for specific types of control design not considered here such as boundary control.

\bibliographystyle{IEEEtran} % Use the plain bibliography style
\bibliography{name}  % The name of your .bib file (references.bib)

\section{Appendix}
\subsection{Additional Examples}\
\begin{example} (More Beam Plots)
Additional plots showing the evolution of actual PDE states for the Timoshenko Beam example are given in Fig. \ref{fig:placeholder1} and \ref{fig:placeholder2}.
\begin{figure}
    \centering
    \includegraphics[width=0.99\linewidth]{newestfigures/ConsistentFigs/dotw.pdf}
    \caption{Evolution of the first PDE state, $\dot{w}$, for the Timoshenko Beam Equations under state-feedback control.}
    \label{fig:placeholder1}
\end{figure}

\begin{figure}
    \centering
    \includegraphics[width=0.99\linewidth]{newestfigures/ConsistentFigs/dotphi.pdf}
    \caption{Evolution of the first PDE state, $\dot{\phi}$, for the Timoshenko Beam Equations under state-feedback control.}
    \label{fig:placeholder2}
\end{figure}
\end{example}
\begin{example}[Transport Equation]
Next, we consider a transport equation with control as
\begin{align*}
    \partial_tx(t,s) &= \partial_{s}x(t,s) + 10s(s-1)(s-0.5)w(t) +u(t),\\
    z(t) &=  \int_0^1 x(t,s)\ ds,\quad x(t, 1) = 0.
\end{align*}
Although the transport equation is naturally stable, the control action can be utilized to suppress vibrations. Similar to unstable reaction-diffusion PDE, we search for an optimal controller using LPIs in Cor. \ref{cor:control} by bisecting on the $\gamma$ parameter. As seen in Fig. \ref{fig:transportresponse}, we see the stability is maintained. Moreover, the uncontrolled system has an impulse-to-peak bound of $0.1678$, whereas with control we can achieve an impulse-to-peak norm of $0.1247$ (see Fig. \ref{fig:controlledregulatedoutputstransport}). The transport problem had a mean iteration time in CPU seconds of 3.43 over 15 iterations with a problem size involing 85271 decision variables and 699 constraints. %the peak output is reduced. In the controlled case, the absolute value of the peak of the regulated output is below 0.02, whereas in the uncontrolled case, the peak has an absolute value between 0.14 and 0.16.
\end{example}



\begin{figure}
    \centering
    \includegraphics[width=0.99\linewidth]{newestfigures/ConsistentFigs/TransportSurf.pdf}
    \caption{Evolution of the PDE state for the transport equation under state-feedback controller}
    \label{fig:transportresponse}
\end{figure}
\begin{figure}
    \centering
    \includegraphics[width=0.8\linewidth]{newestfigures/ConsistentFigs/TransportReg.pdf}
    \caption{The regulated output, $\int_0^1 x(t,s)ds$, for the controlled ($\textcolor{red}{--}$) and uncontrolled ($\textcolor{blue}{-}$) transport equation. The control reduces the peak response under impulsive input in comparison to uncontrolled system.}
    \label{fig:controlledregulatedoutputstransport}
\end{figure}

\end{document}
